% ============================================================
% SI APPENDIX (RESEARCH ARTICLE VERSION)
% "Additive genetic channels: a dynamical theory of when
%  predictive regimes hold in evolution and breeding"
% Ortiz-Barrientos & Cooper — PNAS Research Article
%
% PNAS SI rules (Research Articles):
%   Supporting text, figures, and tables all permitted.
%   Submit as a single separate PDF on PNAScentral.
%   Number SI figures S1, S2, ...; SI tables ST1, ST2, ...
%   Number SI equations (S1), (S2), ...
%
% KEY CHANGES FROM BRIEF REPORT SI:
%   - Brief Report SI limited to extended methods only (no figures).
%   - Research Article SI has no such restriction.
%   - Added: SI Figures S1 (parameter sensitivity), S2 (two-trait
%     nullclines), SI Table ST1 (parameter definitions).
%   - Sections renumbered S1–S7 to match main-text cross-references.
%   - Prose introductions added to each section to aid the
%     non-specialist reader (audience: same as main text).
% ============================================================

\documentclass[11pt]{article}

\usepackage[margin=1in]{geometry}
\usepackage[T1]{fontenc}
\usepackage[utf8]{inputenc}
\usepackage{lmodern}
\usepackage{microtype}
\usepackage{amsmath,amssymb,amsthm}
\usepackage{bm}
\usepackage{booktabs,array,tabularx}
\usepackage{graphicx}
\graphicspath{{./}}
\usepackage[numbers,sort&compress]{natbib}
\usepackage{hyperref}
\usepackage{caption}

\hypersetup{colorlinks=true,citecolor=black,linkcolor=black,
            urlcolor=blue}

% ----------------------------------------------------------
% SI equation numbering: (S1), (S2), ...
% ----------------------------------------------------------
\renewcommand{\theequation}{S\arabic{equation}}

% ----------------------------------------------------------
% SI figure numbering: Fig. S1, S2, ...
% ----------------------------------------------------------
\renewcommand{\thefigure}{S\arabic{figure}}
\renewcommand{\thetable}{S\arabic{table}}

% ----------------------------------------------------------
% Custom commands (must match main text)
% ----------------------------------------------------------
\newcommand{\vect}[1]{\boldsymbol{#1}}
\newcommand{\mat}[1]{\mathbf{#1}}
\newcommand{\Ag}{\mathcal{A}_{g}}
\newcommand{\CalL}{\mathcal{L}}
\newcommand{\Vlin}{V_{\mathrm{lin}}}
\newcommand{\Vquad}{V_{\mathrm{quad}}}
\newcommand{\Gsel}{\mat{G}^{\mathrm{sel}}}
\newcommand{\tr}{\operatorname{tr}}
\newcommand{\Var}{\operatorname{Var}}
\newcommand{\Cov}{\operatorname{Cov}}
\newcommand{\Corr}{\operatorname{Corr}}

\setlength{\parskip}{0.5em}
\setlength{\parindent}{0pt}

% ----------------------------------------------------------
\begin{document}
% ----------------------------------------------------------

\begin{center}
  {\Large\bfseries SI Appendix}\\[0.6em]
  {\large Additive genetic channels: a dynamical theory of when
   predictive regimes hold in evolution and breeding}\\[0.4em]
  {\normalsize Daniel Ortiz-Barrientos and Mark Cooper}
\end{center}

\medskip

\noindent
This SI Appendix contains the full mathematical derivations
underlying the main-text results, together with supporting figures
and parameter tables.
The sections correspond to the order in which results appear in
the main text.
Each section opens with a brief orientation paragraph written for
the non-specialist reader; the formal derivations follow.

\bigskip

\noindent\textbf{Contents}

\begin{enumerate}
  \item Variance decomposition and the additivity index
  \item The logit identity
  \item Transport laws for the local gradient and curvature
  \item The closed state-space system
  \item Mechanistic decomposition of $r(t)$
  \item One-trait moving-optimum model
  \item Introgression criterion
\end{enumerate}

Supporting figures: S1 (parameter sensitivity of $r(t)$),
S2 (two-trait state-space nullclines).

Supporting table: ST1 (symbol definitions).

% ==========================================================
\section*{S1.~Variance decomposition and the additivity index}
\label{sec:decomp}
% ==========================================================

\noindent\textit{Orientation.}
This section justifies the variance decomposition in
Eq.~(2) of the main text.
The key steps are: (i) write log-fitness as a linear term plus a
quadratic correction; (ii) show that under Gaussian breeding-value
distributions their variances add without a cross-term;
(iii) compute each component from the gradient, Hessian, and
covariance matrix; and (iv) show that the additivity index equals
the squared correlation between the linear predictor and
log-fitness (a local $R^2$).
The derivation also recovers Lande's response equation and
Fisher's first-order fitness increase as special cases, confirming
that the framework nests existing theory.

\subsection*{S1.1\quad Setup}

All coordinates are breeding-value coordinates.
Let $\vect{z}\in\mathbb{R}^n$ denote breeding value,
$\bar{\vect{z}}$ the population mean, and
$\mat{G} = \Cov(\vect{z})$ the additive genetic covariance matrix.
The second-order local log-fitness expansion is
\begin{equation}
\ell(\vect{z})
\approx \ell(\bar{\vect{z}})
+ \vect{b}^\top\!\Delta\vect{z}
+ \tfrac{1}{2}\Delta\vect{z}^\top\mat{H}\Delta\vect{z},
\label{eq:S1_taylor}
\end{equation}
where $\Delta\vect{z} = \vect{z} - \bar{\vect{z}}$,
$\vect{b} = \nabla\ell(\bar{\vect{z}})$ is the local gradient,
and $\mat{H} = \nabla^2\ell(\bar{\vect{z}})$ is the local Hessian.
The Hessian is negative semidefinite for a fitness peak.
Define the linear and quadratic parts of the log-fitness deviation:
\begin{equation}
L = \vect{b}^\top\Delta\vect{z},
\qquad
Q = \tfrac{1}{2}\Delta\vect{z}^\top\mat{H}\Delta\vect{z},
\label{eq:S1_LQ}
\end{equation}
so $\delta\ell \approx L + Q$.

\subsection*{S1.2\quad Gaussian closure and vanishing of the
             cross-term}

Assume $\Delta\vect{z} \sim \mathcal{N}(\vect{0},\mat{G})$.
The cross-covariance is
$\Cov(L,Q) = \mathbb{E}[LQ] - \mathbb{E}[L]\,\mathbb{E}[Q]$.
Since $\mathbb{E}[L] = \vect{b}^\top\mathbb{E}[\Delta\vect{z}] = 0$,
it suffices to show $\mathbb{E}[LQ] = 0$:
\begin{equation}
\mathbb{E}[LQ]
= \tfrac{1}{2}\sum_{i,j,k}
  b_i H_{jk}\,
  \mathbb{E}[\Delta z_i\,\Delta z_j\,\Delta z_k] = 0,
\label{eq:S1_cross}
\end{equation}
because all odd central moments of a centred multivariate Gaussian
vanish: $\mathbb{E}[\Delta z_i\,\Delta z_j\,\Delta z_k] = 0$.
This is a parity argument.
$L$ is odd under
$\Delta\vect{z}\mapsto-\Delta\vect{z}$; $Q$ is even.
A symmetric distribution integrates odd $\times$ even to zero.

\subsection*{S1.3\quad Variance components}

Since $\Cov(L,Q) = 0$:
\begin{equation}
\Var[\delta\ell] = \Vlin + \Vquad,
\label{eq:S1_decomp}
\end{equation}
with
\begin{align}
\Vlin &= \Var[L] = \vect{b}^\top\mat{G}\vect{b},
\label{eq:S1_Vlin} \\[4pt]
\Vquad &= \Var[Q] = \tfrac{1}{2}\tr[(\mat{H}\mat{G})^2].
\label{eq:S1_Vquad}
\end{align}
Equation~\eqref{eq:S1_Vlin} follows from
$\Var[\vect{b}^\top\Delta\vect{z}]
= \vect{b}^\top\mat{G}\vect{b}$ and requires only the existence of
$\mat{G}$; Gaussianity is not needed here.
Equation~\eqref{eq:S1_Vquad} follows from Isserlis' (Wick's)
theorem for fourth moments of a centred Gaussian:
$\mathbb{E}[\Delta z_i\Delta z_j\Delta z_k\Delta z_l]
= G_{ij}G_{kl} + G_{ik}G_{jl} + G_{il}G_{jk}$,
which contracts to $\tfrac{1}{2}\tr[(\mat{H}\mat{G})^2]$ after
index summation.
One copy of $\mat{G}$ appears in $\Vlin$ and two copies appear in
$\Vquad$.
Under uniform rescaling $\mat{G}\mapsto s\mat{G}$ with fixed
$\vect{b}$ and $\mat{H}$:
\begin{equation}
\Vlin(s) = s\,\Vlin, \qquad \Vquad(s) = s^2\,\Vquad.
\label{eq:S1_scaling}
\end{equation}
Variance compression therefore deepens a channel because
$\Vquad$ falls faster than $\Vlin$.

\subsection*{S1.4\quad The additivity index as a local $R^2$}

The additivity index $\Ag = \Vlin/(\Vlin+\Vquad)$.
Because $\delta\ell = L + Q$ and $\Cov(L,Q) = 0$:
\begin{equation}
\Corr(L,\,\delta\ell)^2
= \frac{\Var[L]}{\Var[L]+\Var[Q]}
= \frac{\Vlin}{\Vlin+\Vquad}
= \Ag.
\label{eq:S1_R2}
\end{equation}
$\Ag$ is therefore the coefficient of determination of the best
linear predictor of log-fitness from breeding values.

\subsection*{S1.5\quad Recovering Lande's equation and Fisher's
             theorem}

Under weak selection, the per-generation mean change is
$\Delta\bar{\vect{z}} \approx \Cov[\Delta\vect{z},\,\delta\ell]$.
The same parity argument gives
$\Cov[\Delta\vect{z},Q] = \vect{0}$, so only $L$ contributes:
\begin{equation}
\Delta\bar{\vect{z}} = \mat{G}\vect{b},
\label{eq:S1_lande}
\end{equation}
which is Lande's multivariate response equation \cite{Lande1979}.
Premultiplying by $\vect{b}^\top$:
\begin{equation}
\Delta\bar{\ell}
\approx \vect{b}^\top\mat{G}\vect{b}
= \Vlin,
\label{eq:S1_fisher}
\end{equation}
which is Fisher's first-order fitness increase.
Both results rest on the same Gaussian closure as
Eq.~\eqref{eq:S1_decomp}: they are local tangent-plane statements
valid precisely when $\Ag$ is high.

% ==========================================================
\section*{S2.~The logit identity}
\label{sec:logit}
% ==========================================================

\noindent\textit{Orientation.}
The additivity index $\Ag$ is bounded between 0 and 1, which
makes rates of change awkward to interpret near the boundaries.
A logit transform maps it to an unbounded scale and produces a
clean relationship between the rate diagnostic $r(t)$ and the
rate of change of $\Ag$.
This section derives that relationship.
The key result is Eq.~(S12): $\dot{\Ag} = \Ag(1-\Ag)\,r(t)$.
This is a pure algebraic identity --- it follows from the
definition of $r(t)$ alone, without any additional biological
assumptions.

Define the logit transform
\begin{equation}
\CalL(t) = \ln\!\left(\frac{\Vlin(t)}{\Vquad(t)}\right)
          = \ln\!\left(\frac{\Ag(t)}{1-\Ag(t)}\right).
\label{eq:S2_logit}
\end{equation}
Differentiating:
\begin{equation}
r(t) \equiv \dot{\CalL}
= \frac{\dot{\Vlin}}{\Vlin} - \frac{\dot{\Vquad}}{\Vquad}.
\label{eq:S2_r_def}
\end{equation}
Differentiate $\Ag = \Vlin/(\Vlin+\Vquad)$ with respect to time:
\begin{align}
\dot{\Ag}
&= \frac{\dot{\Vlin}\,\Vquad - \Vlin\,\dot{\Vquad}}
        {(\Vlin+\Vquad)^2}
 = \frac{\Vlin\,\Vquad}{(\Vlin+\Vquad)^2}
   \left(\frac{\dot{\Vlin}}{\Vlin}-\frac{\dot{\Vquad}}{\Vquad}\right)
\notag \\
&= \Ag(1-\Ag)\,r(t).
\label{eq:S2_logistic}
\end{align}
The factor $\Vlin\,\Vquad/(\Vlin+\Vquad)^2 = \Ag(1-\Ag)$
follows from the definitions.
Integrating:
\begin{equation}
\CalL(t) = \CalL(0) + \int_0^t r(s)\,\mathrm{d}s.
\label{eq:S2_integral}
\end{equation}
Re-entry into the channel after exit at time $T$ occurs when
$\int_T^t r(s)\,\mathrm{d}s > 0$.
The cumulative integral of $r(t)$ is therefore the quantity to
monitor for forecasting re-entry.

% ==========================================================
\section*{S3.~Transport laws for the local gradient and curvature}
\label{sec:transport}
% ==========================================================

\noindent\textit{Orientation.}
The gradient $\vect{b}$ and Hessian $\mat{H}$ at the current
population mean are functions of two moving quantities: the
mean itself and the environmental state.
Their rates of change follow from the chain rule alone.
No dynamical model for selection is needed here; the equations
are purely geometric.
Biologically, this section says: if the population moves across
a curved surface, the slope at its feet changes; if the
environment shifts the shape of the surface, the slope changes
at fixed position.
Both effects are captured by the transport laws below.

Let $u$ be a scalar environmental state.
Write $b_i(t) = \partial\ell/\partial z_i(\bar{\vect{z}},u)$ and
differentiate with respect to time:
\begin{equation}
\dot{b}_i
= \sum_j H_{ij}\,\dot{\bar{z}}_j
+ \frac{\partial^2\ell}{\partial z_i\partial u}\,\dot{u}.
\label{eq:S3_bdot_comp}
\end{equation}
In matrix form, defining
$({\vect{b}}_u)_i = \partial^2\ell/\partial z_i\partial u$:
\begin{equation}
\dot{\vect{b}}
= \mat{H}\,\dot{\bar{\vect{z}}} + \vect{b}_u\,\dot{u}.
\label{eq:S3_bdot}
\end{equation}
Similarly, differentiating
$H_{ij}(t) = \partial^2\ell/\partial z_i\partial z_j$:
\begin{equation}
\dot{\mat{H}}
= \mat{T}[\dot{\bar{\vect{z}}}] + \mat{H}_u\,\dot{u},
\label{eq:S3_Hdot}
\end{equation}
where $T_{ijk} = \partial^3\ell/\partial z_i\partial z_j\partial z_k$
is the third-derivative tensor,
$\mat{T}[\vect{v}]_{ij} = \sum_k T_{ijk}v_k$, and
$(\mat{H}_u)_{ij} = \partial^3\ell/\partial z_i\partial z_j\partial u$.
A dynamical theory of curvature necessarily involves third derivatives
of the fitness surface: understanding how the Hessian changes as the
population moves requires knowing how the gradient of the Hessian
is oriented.
This is one derivative order above what is needed to move the mean,
and it is the reason that estimating $\dot{\mat{H}}$ requires
more data than estimating $\dot{\vect{b}}$.

% ==========================================================
\section*{S4.~The closed state-space system}
\label{sec:statespace}
% ==========================================================

\noindent\textit{Orientation.}
This section assembles the transport laws from Section~S3
with the selection step for the mean and covariance, plus a
minimal closure for the transmission processes (recombination,
mutation, introgression, drift), into a single five-equation
dynamical system.
The important point is that this system is closed: given initial
conditions and a fitness surface $\ell(\vect{z},u)$, it evolves
all five state variables without further biological assumptions.
The additivity index $\Ag$ is an output of the system, not a
parameter.

Under Gaussian closure and weak selection, the per-generation
changes driven by selection are (derived in \cite{OrtizBarrientosCooper2026}):
\begin{align}
\dot{\bar{\vect{z}}}_{\mathrm{sel}} &= \mat{G}\vect{b},
\label{eq:S4_zdot} \\[4pt]
\dot{\mat{G}}_{\mathrm{sel}} &= \mat{G}\mat{H}\mat{G}.
\label{eq:S4_Gdot_sel}
\end{align}
Equation~\eqref{eq:S4_zdot} is Lande's response law.
Equation~\eqref{eq:S4_Gdot_sel} is its companion: curvature
reshapes the covariance cloud through the same local weighting
that moves the mean.
Both arise from completing the square in the post-selection
Gaussian kernel.

Adding a minimal transmission closure:
\begin{equation}
\dot{\mat{G}}
= \mat{G}\mat{H}\mat{G}
  - \lambda_{\mathrm{rec}}(\mat{G} - \mat{G}_*)
  + \mat{M}
  + \mat{J}(t)
  - \frac{1}{2N_e}\mat{G},
\label{eq:S4_Gdot}
\end{equation}
where $\mat{G}_*$ is the recombination baseline, $\lambda_{\mathrm{rec}}$
is the recombination relaxation rate, $\mat{M}$ is mutational
input per generation, $\mat{J}(t)$ is introgression or migration
variance, and $N_e$ is effective population size.
Substituting $\dot{\bar{\vect{z}}} = \mat{G}\vect{b}$ into
the transport laws from Section~S3:
\begin{align}
\dot{\vect{b}} &= \mat{H}\mat{G}\vect{b} + \vect{b}_u\,\dot{u},
\label{eq:S4_bdot} \\
\dot{\mat{H}} &= \mat{T}[\mat{G}\vect{b}] + \mat{H}_u\,\dot{u}.
\label{eq:S4_Hdot}
\end{align}
Together with an environmental equation
$\dot{u} = g(u,\bar{\vect{z}},t)$,
Eqs.~\eqref{eq:S4_zdot}, \eqref{eq:S4_Gdot},
\eqref{eq:S4_bdot}, \eqref{eq:S4_Hdot}, and the environmental
equation form the closed state-space system.

See Fig.~S2 for nullcline structure of the two-trait system.

% ==========================================================
\section*{S5.~Mechanistic decomposition of $r(t)$}
\label{sec:decomp_r}
% ==========================================================

\noindent\textit{Orientation.}
This section derives the three-way split
$r = r_{\mathrm{motion}} + r_{\mathrm{env}} + r_{\mathrm{cov}}$
and the five-way split within $r_{\mathrm{cov}}$.
The key insight is that each biological process (motion, environmental
forcing, selection, recombination, mutation, introgression, drift)
contributes a signed component to the rate diagnostic, and those
signs have direct biological interpretations.
The drift term is always positive; the mutation term is always
negative; the others depend on alignment between the relevant
variance change and the fitness gradient.

Differentiate $\Vlin = \vect{b}^\top\mat{G}\vect{b}$ and
$\Vquad = \tfrac{1}{2}\tr[(\mat{H}\mat{G})^2]$:
\begin{align}
\dot{\Vlin}
&= 2\,\vect{b}^\top\mat{G}\dot{\vect{b}}
  + \vect{b}^\top\dot{\mat{G}}\vect{b},
\label{eq:S5_Vlindot} \\[4pt]
\dot{\Vquad}
&= \tr[\mat{G}\mat{H}\mat{G}\,\dot{\mat{H}}]
  + \tr[\mat{H}\mat{G}\mat{H}\,\dot{\mat{G}}].
\label{eq:S5_Vquaddot}
\end{align}
Substituting the closed equations from Section~S4
and collecting by source:
\begin{align}
r_{\mathrm{motion}}
&= \frac{2\vect{b}^\top\mat{G}\mat{H}\mat{G}\vect{b}}{\Vlin}
   - \frac{\tr[\mat{G}\mat{H}\mat{G}\,\mat{T}[\mat{G}\vect{b}]]}
          {\Vquad},
\label{eq:S5_r_motion} \\[6pt]
r_{\mathrm{env}}
&= \frac{2\vect{b}^\top\mat{G}\vect{b}_u\,\dot{u}}{\Vlin}
   - \frac{\tr[\mat{G}\mat{H}\mat{G}\mat{H}_u]\,\dot{u}}{\Vquad},
\label{eq:S5_r_env} \\[6pt]
r_{\mathrm{cov}}
&= \frac{\vect{b}^\top\dot{\mat{G}}\vect{b}}{\Vlin}
   - \frac{\tr[\mat{H}\mat{G}\mat{H}\,\dot{\mat{G}}]}{\Vquad}.
\label{eq:S5_r_cov}
\end{align}
Substituting Eq.~\eqref{eq:S4_Gdot} into
Eq.~\eqref{eq:S5_r_cov} and collecting by transmission process:
\begin{align}
r_{\mathrm{sel}}
&= \frac{\vect{b}^\top\mat{G}\mat{H}\mat{G}\vect{b}}{\Vlin}
   - \frac{\tr[\mat{H}\mat{G}\mat{H}\mat{G}\mat{H}\mat{G}]}
          {\Vquad},
\label{eq:S5_r_sel} \\[6pt]
r_{\mathrm{rec}}
&= -\lambda_{\mathrm{rec}}
   \!\left[
   \frac{\vect{b}^\top(\mat{G}-\mat{G}_*)\vect{b}}{\Vlin}
   - \frac{\tr[\mat{H}\mat{G}\mat{H}(\mat{G}-\mat{G}_*)]}
          {\Vquad}
   \right],
\label{eq:S5_r_rec} \\[6pt]
r_{\mathrm{mut}}
&= \frac{\vect{b}^\top\mat{M}\vect{b}}{\Vlin}
   - \frac{\tr[\mat{H}\mat{G}\mat{H}\mat{M}]}{\Vquad},
\label{eq:S5_r_mut} \\[6pt]
r_{\mathrm{intro}}
&= \frac{\vect{b}^\top\mat{J}(t)\vect{b}}{\Vlin}
   - \frac{\tr[\mat{H}\mat{G}\mat{H}\mat{J}(t)]}{\Vquad},
\label{eq:S5_r_intro} \\[6pt]
r_{\mathrm{drift}}
&= +\frac{1}{2N_e}.
\label{eq:S5_r_drift}
\end{align}
The positive sign of $r_{\mathrm{drift}}$ is a direct consequence
of the scaling asymmetry (Eq.~\ref{eq:S1_scaling}): drift reduces
$\Vlin$ linearly in $\mat{G}$ but $\Vquad$ quadratically, so
local additivity increases mechanically under drift even as
evolvability declines.

See Fig.~S1 for the sensitivity of $r(t)$ to the parameters
$s$, $G$, $d$, and $N_e$ in the one-trait model.

% ==========================================================
\section*{S6.~One-trait moving-optimum model}
\label{sec:onetrait}
% ==========================================================

\noindent\textit{Orientation.}
The one-trait model is the simplest case in which all five
quantities --- $d$, $G$, $b$, $H$, and $r$ --- can be written in
closed form.
It serves as the testbed for every result in the main text.
The model is also the cleanest illustration of the
counterintuitive result that a population near its optimum is
less additively predictable than one far from it, because the
directional signal collapses at the peak.

Consider the scalar fitness surface
\begin{equation}
\ell(z,t)
= \ell_0 - \tfrac{1}{2}s(z - \theta)^2,
\label{eq:S6_landscape}
\end{equation}
where $\theta(t)$ is the moving optimum and $s > 0$ is
stabilizing curvature.
Define the lag $d = \theta - \bar{z}$.
At the current mean:
\begin{equation}
b = s\,d, \qquad H = -s.
\label{eq:S6_bH}
\end{equation}

\subsection*{S6.1\quad Diagnostic layer}

Substituting into the variance formulae:
\begin{equation}
\Vlin = (sd)^2 G = s^2d^2G,
\qquad
\Vquad = \tfrac{1}{2}s^2G^2.
\label{eq:S6_VV}
\end{equation}
The $s^2$ factor cancels in the ratio:
\begin{equation}
\Ag = \frac{d^2G}{d^2G + \tfrac{1}{2}G^2}
    = \frac{d^2}{d^2 + G/2}.
\label{eq:S6_Ag}
\end{equation}
The channel condition $\Ag > \alpha$ translates to
$d^2 > \alpha G/(2(1-\alpha))$:
the squared lag must exceed a fixed multiple of the variance.

\subsection*{S6.2\quad State equations}

The lag equation follows from
$\dot{\bar{z}} = Gb = Gsd$:
\begin{equation}
\dot{d} = \dot{\theta} - Gsd.
\label{eq:S6_ddot}
\end{equation}
The one-dimensional covariance closure:
\begin{equation}
\dot{G}
= -sG^2
  - \lambda_{\mathrm{rec}}(G - G_*)
  + m
  - \frac{G}{2N_e}.
\label{eq:S6_Gdot}
\end{equation}

\subsection*{S6.3\quad Rate diagnostic}

Differentiating $\CalL = \ln(2d^2/G)$ and substituting:
\begin{equation}
r
= 2\frac{\dot{\theta}}{d}
  - sG
  + \lambda_{\mathrm{rec}}\!\left(1 - \frac{G_*}{G}\right)
  - \frac{m}{G}
  + \frac{1}{2N_e}.
\label{eq:S6_r}
\end{equation}
The five terms correspond exactly to the five-way decomposition of
$r_{\mathrm{cov}}$ plus the motion term:
$2\dot{\theta}/d$ regenerates directional signal;
$-sG$ increases curvature exposure;
$\lambda_{\mathrm{rec}}(1 - G_*/G)$ reflects recombination
relative to the transmission baseline;
$-m/G$ reflects isotropic mutation input;
$+1/(2N_e)$ is the drift contribution.

\subsection*{S6.4\quad Distinguishing Act~II from Act~III}

\textbf{Act~II} (optimum reached):
$\dot{\theta} \approx 0$, so $r_{\mathrm{motion}} \to 0$.
The dominant term is $-sG < 0$, driving steady channel exit.
$r_{\mathrm{env}} \approx 0$ throughout.

\textbf{Act~III} (gradient rotation; two-trait model):
The environmental forcing term $r_{\mathrm{env}}$ becomes strongly
negative as $\vect{b}$ rotates away from the major axis of
$\mat{G}$, because $\vect{b}_u$ has large projection onto the
misaligned direction.
Meanwhile, $r_{\mathrm{cov}}$ can remain positive if selection
is still compressing the cloud along the old gradient.
The sign pattern $r_{\mathrm{env}} \ll 0$ with $r_{\mathrm{cov}} > 0$
does not arise in Act~II ($r_{\mathrm{env}} \approx 0$,
$r_{\mathrm{cov}} < 0$).
The two scenarios are therefore distinguishable from the
decomposition of $r(t)$ before prediction accuracy deteriorates.

% ==========================================================
\section*{S7.~Introgression criterion}
\label{sec:introgression}
% ==========================================================

\noindent\textit{Orientation.}
When new genetic variance enters a population through a wide cross
or migration, it can either help or hurt local predictability
depending on whether it is aligned with the fitness gradient or
exposed to curvature.
This section derives the first-order change in $\Ag$ from a rank-1
variance input along direction $\vect{v}$, and identifies the
optimal direction for incoming variance.

Suppose new variance is introduced along direction $\vect{v}$:
\begin{equation}
\mat{G}^+ = \mat{G} + \sigma_\Delta^2\,\vect{v}\vect{v}^\top.
\label{eq:S7_Gplus}
\end{equation}
To first order in $\sigma_\Delta^2$:
\begin{align}
\Vlin^+ &= \Vlin + \sigma_\Delta^2(\vect{b}^\top\vect{v})^2,
\label{eq:S7_Vlin} \\[4pt]
\Vquad^+ &= \Vquad
           + \sigma_\Delta^2\,\vect{v}^\top\mat{H}\mat{G}\mat{H}\vect{v}
           + O(\sigma_\Delta^4).
\label{eq:S7_Vquad}
\end{align}
Eq.~\eqref{eq:S7_Vlin} follows from bilinearity of the quadratic
form; Eq.~\eqref{eq:S7_Vquad} from expanding
$\tr[(\mat{H}\mat{G}^+)^2]$ to first order.
The first-order change in $\Ag$:
\begin{equation}
\delta\Ag
\approx \Ag(1-\Ag)
\left[
  \frac{(\vect{b}^\top\vect{v})^2}{\Vlin}
  -
  \frac{\vect{v}^\top\mat{H}\mat{G}\mat{H}\vect{v}}{\Vquad}
\right]
\sigma_\Delta^2.
\label{eq:S7_deltaAg}
\end{equation}
Introgression deepens the channel ($\delta\Ag > 0$) when
gradient-alignment benefit
$(\vect{b}^\top\vect{v})^2/\Vlin$
exceeds curvature-exposure cost
$\vect{v}^\top\mat{H}\mat{G}\mat{H}\vect{v}/\Vquad$.

The optimal introgression direction maximises $\delta\Ag$ over unit
vectors $\vect{v}$.
This is the leading eigenvector of
$\mat{G}^{-1/2}\vect{b}\vect{b}^\top\mat{G}^{-1/2}
 - (\Vlin/\Vquad)\mat{G}^{1/2}\mat{H}^2\mat{G}^{1/2}$.
When curvature is isotropic ($\mat{H} = -\gamma\mat{I}$),
this reduces to $\mat{G}^{-1/2}\vect{b}$:
new variance should be gradient-aligned in the metric of
$\mat{G}^{-1}$.

% ==========================================================
% SUPPORTING FIGURES
% ==========================================================

% ==========================================================
% SUPPORTING FIGURES
%
% SI FIGURE PLACEMENT RULES:
%   The SI is a standalone single-column article (11pt).
%   All figures use the plain \begin{figure}[ht] environment.
%   Width recommendations:
%     \linewidth       : full text width — use for multi-panel composites
%     0.80\linewidth   : slightly inset — use for single-panel figures
%     0.65\linewidth   : narrow inset — use for small square figures
%   All filenames are WITHOUT extension; pdflatex finds .pdf files.
% ==========================================================

% ----------------------------------------------------------
% SI FIGURE S1  — parameter sensitivity (2×2 panels)
% 10.0 × 7.5 inch figure; use \linewidth for full text width.
% ----------------------------------------------------------
\begin{figure}[ht]
\centering
\includegraphics[width=\linewidth]{figS1_parameter_sensitivity}
\caption{\textbf{Parameter sensitivity of $r(t)$ in the one-trait
  moving-optimum model.}
  Each panel varies one parameter while holding others at the
  baseline ($s = 0.10$, $G = 1.0$, $d = 2.0$, $N_e = 500$,
  $\lambda_{\mathrm{rec}} = 0.30$, $G_* = 1.0$, $m = 0.005$,
  $\dot\theta = 0.02$).
  Green shading: $r > 0$ (channel improving or deepening).
  Orange shading: $r < 0$ (channel eroding or exiting).
  The red dashed vertical line marks the baseline value.
  \textbf{(A)}~Varying stabilising curvature $s$: stronger
  curvature ($-sG$ term) accelerates channel erosion.
  \textbf{(B)}~Varying lag $d$: large lags sustain the
  directional signal ($2\dot\theta/d$) and keep $r > 0$;
  small lags expose the population to the dominance of
  the $-sG$ term.
  \textbf{(C)}~Varying genetic variance $G$: higher variance
  increases curvature exposure ($-sG$) faster than gradient
  alignment ($\propto G$), so $r$ decreases with $G$.
  \textbf{(D)}~Varying effective population size $N_e$: the
  drift term $+1/(2N_e)$ is always positive but small; smaller
  populations are locally more additive by the scaling
  asymmetry, at the cost of reduced evolvability.}
\label{fig:S1}
\end{figure}

% ----------------------------------------------------------
% SI FIGURE S2  — state-space slice in (s, kappa) plane
% 6.8 × 5.5 inch figure; 0.80\linewidth gives comfortable margins.
% This is Panel B from the original Figure 2, moved to SI.
% ----------------------------------------------------------
\begin{figure}[ht]
\centering
\includegraphics[width=0.80\linewidth]{figS2_statespace_slice}
\caption{\textbf{Channel boundary in the $(s,\,\kappa)$
  state-space plane.}
  The additivity index $\Ag(s,\kappa) = 1/(1+s\kappa)$ is
  evaluated on a grid of variance scale $s$ and
  curvature-to-gradient ratio $\kappa$.
  Blue region: interior channel ($\Ag > \alpha + 0.08$).
  Amber region: boundary zone ($|\Ag - \alpha| \leq 0.08$).
  Red region: exterior ($\Ag < \alpha - 0.08$).
  The red dashed hyperbola is the channel boundary
  $s = (1-\alpha)/(\alpha\kappa)$.
  Three trajectories illustrate qualitatively distinct dynamics.
  $T_1$ (green): genetic variance compresses under directional
  selection, moving the population left in $s$ and deeper
  into the channel.
  $T_2$ (orange): both $s$ and $\kappa$ rise — a boundary
  approach trajectory that does not cross.
  $T_3$ (red): channel exit — the population crosses the
  boundary at the star marker $T_{\mathrm{exit}}$.
  This figure motivates the interpretation of the rate
  diagnostic $r(t)$: trajectories with $r > 0$ move left
  (toward smaller $s$) or down (toward smaller $\kappa$),
  both of which deepen the channel.}
\label{fig:S2}
\end{figure}

% ==========================================================
% SUPPORTING TABLE
% ==========================================================

\begin{table}[ht]
\caption{\textbf{Symbol definitions used throughout the paper and
  SI Appendix.}}
\label{tab:S1}
\small
\begin{tabularx}{\textwidth}{lX}
\toprule
Symbol & Definition \\
\midrule
$\vect{z}$           & Breeding-value vector ($n$-dimensional) \\
$\bar{\vect{z}}$     & Population mean breeding value \\
$\mat{G}$            & Additive genetic covariance matrix \\
$\ell(\vect{z})$     & Local log-fitness surface \\
$\vect{b}$           & Local log-fitness gradient
                       ($= \nabla\ell(\bar{\vect{z}})$) \\
$\mat{H}$            & Local log-fitness Hessian
                       ($= \nabla^2\ell(\bar{\vect{z}})$) \\
$L$                  & Linear part of log-fitness deviation
                       ($= \vect{b}^\top\Delta\vect{z}$) \\
$Q$                  & Quadratic part of log-fitness deviation
                       ($= \tfrac{1}{2}\Delta\vect{z}^\top\mat{H}\Delta\vect{z}$) \\
$\Vlin$              & Gradient-aligned log-fitness variance
                       ($= \vect{b}^\top\mat{G}\vect{b}$) \\
$\Vquad$             & Curvature-exposed log-fitness variance
                       ($= \tfrac{1}{2}\tr[(\mat{H}\mat{G})^2]$) \\
$\Ag$                & Additivity index
                       ($= \Vlin/(\Vlin+\Vquad)$) \\
$\alpha$             & Channel threshold (user-specified; $0 < \alpha < 1$) \\
$\CalL$              & Logit of $\Ag$
                       ($= \ln[\Vlin/\Vquad]$) \\
$r(t)$               & Rate diagnostic ($= \dot{\CalL}$) \\
$u$                  & Scalar environmental state \\
$\vect{b}_u$, $\mat{H}_u$ & Mixed derivatives of $\ell$ with respect
                       to environment \\
$\mat{T}[\,\cdot\,]$ & Third-derivative tensor contraction \\
$\lambda_{\mathrm{rec}}$ & Recombination relaxation rate \\
$\mat{G}_*$          & Recombination baseline covariance \\
$\mat{M}$            & Mutational variance input per generation \\
$\mat{J}(t)$         & Introgression or migration variance input \\
$N_e$                & Effective population size \\
$d$                  & Scalar lag ($= \theta - \bar{z}$) \\
$s$                  & Stabilizing curvature strength \\
$\theta(t)$          & Moving optimum (scalar model) \\
$m$                  & Scalar mutational variance input \\
$\vect{g}_{\max}$    & Leading eigenvector of $\mat{G}$
                       ($= e_1(\mat{G})$) \\
$\lambda_1, \lambda_2$ & Eigenvalues of $\mat{G}$
                       ($\lambda_1 \geq \lambda_2$) \\
\bottomrule
\end{tabularx}
\end{table}

% ==========================================================
% REFERENCES
% ==========================================================

\bibliographystyle{pnas-new}
\begin{thebibliography}{20}

\bibitem{OrtizBarrientosCooper2026}
Ortiz-Barrientos D, Cooper M (2026)
Additive channels in curved fitness landscapes.
\textit{bioRxiv} 2026.03.21.713332.

\bibitem{Lande1979}
Lande R (1979)
Quantitative genetic analysis of multivariate evolution, applied
to brain:body size allometry.
\textit{Evolution} 33:402--416.

\bibitem{LandeArnold1983}
Lande R, Arnold SJ (1983)
The measurement of selection on correlated characters.
\textit{Evolution} 37:1210--1226.

\bibitem{PhillipsArnold1989}
Phillips PC, Arnold SJ (1989)
Visualizing multivariate selection.
\textit{Evolution} 43:1209--1222.

\bibitem{Fisher1930}
Fisher RA (1930)
\textit{The Genetical Theory of Natural Selection.}
Clarendon Press, Oxford.

\bibitem{Price1972}
Price GR (1972)
Fisher's `fundamental theorem' made clear.
\textit{Annals of Human Genetics} 36:129--140.

\bibitem{Schluter1996}
Schluter D (1996)
Adaptive radiation along genetic lines of least resistance.
\textit{Evolution} 50:1766--1774.

\end{thebibliography}

\end{document}
